\section{Intersections in a cover and its mirror}
\label{sec:mirrors}
The previous section studies a possible symmetry of one chamber geometry. For regular covers the same issue is expressed at the level of defining relations: the relations must be invariant under a prescribed reflection. We first separate the intersection calculation from that geometric interpretation.

\begin{theorem}[Intersection transport]\label{thm:transport}
Let $\pi:H\to G$ be a homomorphism and let $C\leq A\cap B$ be subgroups of $H$. If $\pi(A)\cap\pi(B)=\pi(C)$, then
\[
 A\cap B=C\quad\Longleftrightarrow\quad
                    A\cap B\cap\ker\pi\leq C.
\]
No finiteness, centrality or surjectivity hypothesis is required.
\end{theorem}
\begin{proof}
The forward implication is immediate. Given $x\in A\cap B$, the image hypothesis supplies $c\in C$ with $\pi(c)=\pi(x)$. Then $xc^{-1}\in A\cap B\cap\ker\pi$. The condition on this kernel puts $xc^{-1}$ in $C$, and therefore puts $x$ in $C$. The other inclusion was a hypothesis.
\end{proof}

\begin{theorem}[The invariant relation kernel]\label{thm:mirror}
Let $U$ be a group, let $\mu\in\Aut(U)$ satisfy $\mu^2=1$, and let $N\normal U$. Set
\[
 \overline N=\mu(N),\quad \Gamma=U/N,\quad
 \overline\Gamma=U/\overline N,\quad
 \Lambda=\im\bigl(U\longrightarrow\Gamma\times\overline\Gamma\bigr).
\]
Then $\Lambda\cong U/(N\cap\overline N)$. The intersection $N\cap\overline N$ is the largest $\mu$-invariant normal subgroup contained in $N$. Thus every quotient $U/J$ with $J\leq N$ and $\mu(J)=J$ maps to $\Lambda$ over $\Gamma$.

The product $N\overline N$ is the smallest $\mu$-invariant normal subgroup containing $N$. Its image $X=(N\overline N)/N$ in $\Gamma$ satisfies
\begin{equation}\label{eq:chiralitykernel}
                X=\{g\in\Gamma:(g,1)\in\Lambda\}.
\end{equation}
In particular $X=1$ if and only if $N=\overline N$.
\end{theorem}
\begin{proof}
The diagonal map has kernel $N\cap\overline N$, giving the first isomorphism. Applying $\mu$ exchanges the two factors of this intersection. If $J\leq N$ and $\mu(J)=J$, applying $\mu$ to its inclusion gives $J\leq\overline N$, proving maximality and the factorization. The two normal subgroups $N,\overline N$ have normal product; applying $\mu$ exchanges its factors, and normality makes their two products equal. Every invariant normal subgroup containing $N$ also contains $\overline N$, hence their product.

A pair $(g,1)$ in the diagonal image is represented by some $u\in\overline N$. Its first coordinate ranges exactly over the image of $\overline N$ in $U/N$, which is $(N\overline N)/N$. Finally $X=1$ gives $\overline N\leq N$. Applying $\mu$ gives the reverse inclusion. The converse is immediate.
\end{proof}

\subsection{The geometric input for string rotation groups}
We state the geometric correspondence used in the cover application rather than hide it in the word ``group''. Let
\[
 W_n=\langle r_0,\ldots,r_{n-1}\mid r_i^2=1,
       [r_i,r_j]=1\ (|i-j|>1)\rangle,
 \qquad U_n=W_n^+,
 \qquad s_i=r_{i-1}r_i.
\]
Here $W_n^+$ is the subgroup of words of even length and $[x,y]=x^{-1}y^{-1}xy$. The generators of $U_n$ have the presentation
\[
 U_n=\langle s_1,\ldots,s_{n-1}\mid(s_i\cdots s_j)^2=1
                                      \ (i<j)\rangle.
\]
For a marked quotient $\Gamma=\langle\sigma_1,\ldots,\sigma_{n-1}\rangle$, put
\[
 \Gamma_{n-1}=\langle\sigma_1,\ldots,\sigma_{n-2}\rangle,
 \quad\Gamma_0=\langle\sigma_2,\ldots,\sigma_{n-1}\rangle,
 \quad\Gamma_{0,n-1}=\langle\sigma_2,\ldots,\sigma_{n-2}\rangle.
\]
The standard string $C^+$ correspondence identifies the rotation groups of chiral or directly regular abstract polytopes with these marked quotients satisfying the intersection property. In its recursive form, for rank at least three the two end groups must themselves satisfy the property and must intersect exactly in the medial group. The rank-two group is a nontrivial cyclic group, finite or infinite; the rank-one rotation group is trivial. This is the correspondence of Schulte--Weiss, with the recursive form cited by Cunningham as Definition 2.1. We use the mathematically consistent rank-two base, correcting the rank-one line printed in the checked author edition.

The mirror is induced by conjugation by $r_0$:
\begin{equation}\label{eq:mirrorrotation}
 \mu(s_1)=s_1^{-1},\qquad \mu(s_2)=s_1^2s_2,
 \qquad\mu(s_i)=s_i\quad(i>2).
\end{equation}
Direct regularity is equivalent to invariance of the relation kernel under this mirror, or equivalently to the corresponding automorphism of the marked quotient. Marked epimorphisms give covers. Without the intersection property, the pertinent quotients describe oriented maniplexes rather than necessarily polytopes. These correspondence statements are retained geometric theorems, not conclusions of Theorem~\ref{thm:transport}.

For an actual chiral polytope with group $U_n/N$, Theorem~\ref{thm:mirror} constructs the smallest regular maniplex covering it: the rotation group is $\Lambda$. Indeed the intersection kernel is mirror invariant, and every regular cover has an invariant kernel contained in $N$. All regular polytope covers of a chiral polytope are included. To see the required orientability, color the flags of the chiral polytope by its two automorphism orbits and pull that coloring back along the covering map; every adjacency changes color. Regularity then identifies this bipartition with the even-word subgroup. There is no assumption of a finite flag set.

\subsection{Local mirrors and embedded kernels}
Let $A,B,C\leq\Lambda$ be generated by the paired first $n-2$, last $n-2$, and middle $n-3$ rotation generators. Let $K,L,M$ denote the facet, vertex-figure and medial section of the original polytope, with compatible base flags. In the first factor $\Gamma$, the three section groups are actual subgroups, not just isomorphism types. Write $X(K),X(L),X(M)$ for their chirality groups.

\begin{lemma}[Local mirror compatibility]\label{lem:localmirror}
The groups $A$ and $B$ are the mirror mixes of the compatibly marked end sections, up to conjugating their second coordinates. For the second projection $\pi:\Lambda\to\overline\Gamma$, their kernels and the medial kernel are
\begin{equation}\label{eq:localX}
 A\cap\ker\pi=X(K)\times1,\quad
 B\cap\ker\pi=X(L)\times1,\quad
 C\cap\ker\pi=X(M)\times1.
\end{equation}
\end{lemma}
\begin{proof}
Identify $\overline\Gamma$ with $\Gamma$ using the isomorphism induced by $\mu$. In these coordinates the paired generators are $(\sigma_i,\mu(s_i)N)$. On the universal facet subgroup, conjugation by $r_0$ is its own intrinsic mirror. Formula~\eqref{eq:chiralitykernel} immediately gives the assertion for $A$.

The intrinsic mirror on the universal vertex subgroup $\langle s_2,\ldots,s_{n-1}\rangle$ is conjugation by $r_1$. Since $r_0=s_1r_1$, conjugation by $r_0$ is conjugation by $r_1$ followed by conjugation by $s_1$. Conjugating the second coordinate by $\sigma_1^{-1}$ therefore turns the global mirror marking into the intrinsic local mirror marking. This changes neither the first coordinate nor whether the second coordinate is the identity. The same calculation applies to the medial universal subgroup $\langle s_2,\ldots,s_{n-2}\rangle$. Applying \eqref{eq:chiralitykernel} to these two local marked groups proves the remaining identities. This argument does not assert that a chiral section has a mirror automorphism of its own group.
\end{proof}

\begin{lemma}[Prime degree and minimality]\label{lem:primeminimal}
Suppose $R$ is the smallest regular maniplex covering a polytope $P$, and $R$ is not a polytope. If a finite regular polytope $Q$ covers $R$ with prime degree, then $Q$ is a minimal regular polytope cover of $P$. If $R$ is itself a polytope, it is the unique minimal regular polytope cover of $P$, without a finiteness assumption.
\end{lemma}
\begin{proof}
Every regular polytope cover of $P$ factors through $R$ by smallestness. An intermediate regular cover under $Q$ therefore corresponds to an intermediate normal subgroup of the prime-order kernel of $\Gamma(Q)\to\Gamma(R)$ (or of the rotation groups in the orientable case). It is either trivial or the whole kernel. The second possibility gives $R$, which is not polytopal, so the only polytope under $Q$ in this order is $Q$ itself. If $R$ is polytopal, it lies under every regular polytope cover of $P$ and is consequently the unique minimal one. This latter argument uses no group orders.
\end{proof}
